% GENERATED FILE. Edit canonical lessons, then run npm run generate:books.
% canonical-source-hash: fnv1a64:daf856060acd9477
% canonical-lessons: KA-C173-adugemane, KA-C173-snanada-mane, KA-C173-saucalaya, KA-C173-kone, KA-C173-angala

\chapter{Kitchen, Bathroom, Toilet, Room, Courtyard}
\label{ch:ka-kitchen-bathroom-toilet-room-courtyard}
\hlchaptermodality{173}

\begin{chapteropening}
\textbf{By the end of this chapter:} \emph{I can say a kitchen, a bathroom, a toilet, a room and a courtyard.}

Complete the last lesson of chapter 173: I can say a kitchen, a bathroom, a toilet, a room and a courtyard.
\end{chapteropening}

\section[\texorpdfstring{\kn{ಅಡುಗೆಮನೆ} (aḍugemane)}{ಅಡುಗೆಮನೆ (aḍugemane)}]{\kn{ಅಡುಗೆಮನೆ} (aḍugemane) — a kitchen}

\label{lesson:KA-C173-adugemane}



\begin{quote}
Before the new one: say the Kannada for a century, then the Kannada for a moment.
\end{quote}

\subsection*{You'll want to know: \kn{ಅಡುಗೆಮನೆ}}
\textbf{\kn{ಅಡುಗೆಮನೆ}} — \emph{aḍugemane} — \textquotedblleft{}a kitchen\textquotedblright{}.

\begin{grammarlens}[title={What you've built}]
One more word for this chapter.
\end{grammarlens}

\subsection*{Your turn}
\begin{itemize}
\raggedright
  \item \emph{Say it:} \emph{aḍugemane}
  \item \emph{Say it:} \emph{aḍugemane}, once more
  \item \emph{Say it:} say \emph{aḍugemane}
  \item \emph{Recall:} say the Kannada for a century, then the Kannada for a moment, then say \emph{aḍugemane} again
\end{itemize}

\begin{tcolorbox}[breakable,colback=teal!4,colframe=teal!35!black,title={Before you move on}]
What does \kn{ಅಡುಗೆಮನೆ} mean? (\textquotedblleft{}A kitchen\textquotedblright{}.) Say it once more.
\end{tcolorbox}
\section[\texorpdfstring{\kn{ಸ್ನಾನದ} \kn{ಮನೆ} (snānada mane)}{ಸ್ನಾನದ ಮನೆ (snānada mane)}]{\kn{ಸ್ನಾನದ} \kn{ಮನೆ} (snānada mane) — a bathroom}

\label{lesson:KA-C173-snanada-mane}



\begin{quote}
Before the new one: say the Kannada for a moment, then the Kannada for a kitchen.
\end{quote}

\subsection*{You'll want to know: \kn{ಸ್ನಾನದ} \kn{ಮನೆ}}
\textbf{\kn{ಸ್ನಾನದ} \kn{ಮನೆ}} — \emph{snānada mane} — \textquotedblleft{}a bathroom\textquotedblright{}.

\begin{grammarlens}[title={What you've built}]
One more word for this chapter.
\end{grammarlens}

\subsection*{Your turn}
\begin{itemize}
\raggedright
  \item \emph{Say it:} \emph{snānada mane}
  \item \emph{Say it:} \emph{snānada mane}, once more
  \item \emph{Say it:} say \emph{snānada mane}
  \item \emph{Recall:} say the Kannada for a moment, then the Kannada for a kitchen, then say \emph{snānada mane} again
\end{itemize}

\begin{tcolorbox}[breakable,colback=teal!4,colframe=teal!35!black,title={Before you move on}]
What does \kn{ಸ್ನಾನದ} \kn{ಮನೆ} mean? (\textquotedblleft{}A bathroom\textquotedblright{}.) Say it once more.
\end{tcolorbox}
\section[\texorpdfstring{\kn{ಶೌಚಾಲಯ} (śaucālaya)}{ಶೌಚಾಲಯ (śaucālaya)}]{\kn{ಶೌಚಾಲಯ} (śaucālaya) — a toilet}

\label{lesson:KA-C173-saucalaya}



\begin{quote}
Before the new one: say the Kannada for a kitchen, then the Kannada for a bathroom.
\end{quote}

\subsection*{You'll want to know: \kn{ಶೌಚಾಲಯ}}
\textbf{\kn{ಶೌಚಾಲಯ}} — \emph{śaucālaya} — \textquotedblleft{}a toilet\textquotedblright{}.

\begin{grammarlens}[title={What you've built}]
One more word for this chapter.
\end{grammarlens}

\subsection*{Your turn}
\begin{itemize}
\raggedright
  \item \emph{Say it:} \emph{śaucālaya}
  \item \emph{Say it:} \emph{śaucālaya}, once more
  \item \emph{Say it:} say \emph{śaucālaya}
  \item \emph{Recall:} say the Kannada for a kitchen, then the Kannada for a bathroom, then say \emph{śaucālaya} again
\end{itemize}

\begin{tcolorbox}[breakable,colback=teal!4,colframe=teal!35!black,title={Before you move on}]
What does \kn{ಶೌಚಾಲಯ} mean? (\textquotedblleft{}A toilet\textquotedblright{}.) Say it once more.
\end{tcolorbox}
\section[\texorpdfstring{\kn{ಕೋಣೆ} (kōṇe)}{ಕೋಣೆ (kōṇe)}]{\kn{ಕೋಣೆ} (kōṇe) — a room}

\label{lesson:KA-C173-kone}



\begin{quote}
Before the new one: say the Kannada for a bathroom, then the Kannada for a toilet.
\end{quote}

\subsection*{You'll want to know: \kn{ಕೋಣೆ}}
\textbf{\kn{ಕೋಣೆ}} — \emph{kōṇe} — \textquotedblleft{}a room\textquotedblright{}.

\begin{grammarlens}[title={What you've built}]
One more word for this chapter.
\end{grammarlens}

\subsection*{Your turn}
\begin{itemize}
\raggedright
  \item \emph{Say it:} \emph{kōṇe}
  \item \emph{Say it:} \emph{kōṇe}, once more
  \item \emph{Say it:} say \emph{kōṇe}
  \item \emph{Recall:} say the Kannada for a bathroom, then the Kannada for a toilet, then say \emph{kōṇe} again
\end{itemize}

\begin{tcolorbox}[breakable,colback=teal!4,colframe=teal!35!black,title={Before you move on}]
What does \kn{ಕೋಣೆ} mean? (\textquotedblleft{}A room\textquotedblright{}.) Say it once more.
\end{tcolorbox}
\section[\texorpdfstring{\kn{ಅಂಗಳ} (aṅgaḷa)}{ಅಂಗಳ (aṅgaḷa)}]{\kn{ಅಂಗಳ} (aṅgaḷa) — a courtyard, a front yard}

\label{lesson:KA-C173-angala}



\begin{quote}
Before the new one: say the Kannada for a toilet, then the Kannada for a room.
\end{quote}

\subsection*{You'll want to know: \kn{ಅಂಗಳ}}
\textbf{\kn{ಅಂಗಳ}} — \emph{aṅgaḷa} — \textquotedblleft{}a courtyard, a front yard\textquotedblright{}.

\begin{grammarlens}[title={What you've built}]
One more word for this chapter.
\end{grammarlens}

\subsection*{Your turn}
\begin{itemize}
\raggedright
  \item \emph{Say it:} \emph{aṅgaḷa}
  \item \emph{Say it:} \emph{aṅgaḷa}, once more
  \item \emph{Say it:} say \emph{aṅgaḷa}
  \item \emph{Recall:} say the Kannada for a toilet, then the Kannada for a room, then say \emph{aṅgaḷa} again
\end{itemize}

\begin{tcolorbox}[breakable,colback=teal!4,colframe=teal!35!black,title={Before you move on}]
What does \kn{ಅಂಗಳ} mean? (\textquotedblleft{}A courtyard, a front yard\textquotedblright{}.) Say it once more.
\end{tcolorbox}
